% chapters/13/evaluasi.tex
\bagian{Asesmen dan Instrumen Penilaian}
\label{sec:asesmen-level}

\begin{asesmenbox}
\textbf{Kedudukan asesmen bab ini.} Bab ini menuntaskan \SUBCPMK{3.2} dan
menutup \CPMK{3}. Ketiga indikatornya diperiksa melalui penilaian level
(Tugas Praktikum 10) dan soal analisis.

\vspace{0.4em}
Level yang dihasilkan pada bab ini menjadi bagian dari \textbf{Proyek Akhir}
yang dikumpulkan pada minggu kelima belas dan keenam belas.
\end{asesmenbox}

\subsection*{Instrumen 1: Rubrik Praktikum 10}

\begin{table}[htbp]
\centering
\caption{Rubrik penilaian Praktikum 10 beserta indikator yang diukur}
\label{tab:rubrik-praktikum-10}
{\setlength{\tabcolsep}{4pt}\footnotesize%
\begin{tabularx}{\textwidth}{@{}L{3.2cm} C{1.3cm} C{2.1cm} Y@{}}
\toprule
\headrow \textbf{Kriteria} & \textbf{Bobot} & \textbf{Indikator} & \textbf{Deskripsi capaian} \\
\midrule
Penyusunan level & 30\% & \ind{a} &
Level tersusun dari lantai, sekurang-kurangnya tiga bongkahan, dan dua dinding;
seluruh celahnya terjangkau menurut ukuran yang ditetapkan. \\
\addlinespace
Fisika dan kontrol pemain & 40\% & \ind{b} &
Pemain berjalan, melompat, mendarat, dan terhalang dengan benar;
\istilah{Collider} dan \istilah{Rigidbody} terkonfigurasi tepat; daerah
\istilah{trigger} memanggil fungsinya. \\
\addlinespace
Analisis masalah fisika & 30\% & \ind{c} &
Catatan praktikum memuat masalah fisika yang ditemui, gejala, penyebab, dan
perbaikan yang dilakukan, serta hasil percobaan menimbulkan masalah dengan
sengaja. \\
\bottomrule
\end{tabularx}}
\end{table}

\subsection*{Instrumen 2: Soal Analisis}

Jawablah dengan singkat. Kunci jawaban beserta indikator yang diukur disajikan
pada Tabel~\ref{tab:kunci-konsep-ch13}.

\begin{enumerate}
  \item Jelaskan mengapa benda diam tidak perlu diberi \istilah{Rigidbody},
        dan apa akibatnya apabila tetap diberikan.
  \item Sebuah benda bergerak 30 satuan per detik dan lantainya tebal 0,1
        satuan pada langkah fisika 0,02 detik. Tentukan apakah benda itu dapat
        menembus lantai, beserta perhitungannya.
  \item Pemain bergetar ketika berjalan. Sebutkan dua penyebab yang mungkin
        beserta cara memeriksanya.
  \item Sebutkan tiga syarat agar fungsi \istilah{trigger} dipanggil, dan
        jelaskan gejala yang muncul apabila salah satunya tidak dipenuhi.
  \item Mengapa kamera dipindahkan pada \code{LateUpdate}, bukan pada
        \code{Update}?
\end{enumerate}

\begin{table}[htbp]
\centering
\caption{Kunci soal analisis Bab XIII beserta indikator yang diukur}
\label{tab:kunci-konsep-ch13}
{\setlength{\tabcolsep}{4pt}\footnotesize%
\begin{tabularx}{\textwidth}{@{}C{0.8cm} C{2.1cm} Y@{}}
\toprule
\headrow \textbf{No} & \textbf{Indikator} & \textbf{Pokok jawaban} \\
\midrule
1 & \ind{b} &
Karena benda diam tidak bergerak, sehingga tidak ada yang perlu dihitung.
Apabila tetap diberi \istilah{Rigidbody}, mesin menghitungnya tanpa manfaat,
dan benda itu dapat terdorong secara tidak sengaja oleh benda lain. \\
\addlinespace
2 & \ind{c} &
Perpindahan pada satu langkah fisika adalah $30 \times 0{,}02 = 0{,}6$
satuan, yaitu lebih tebal daripada lantainya, sehingga benda dapat menembusnya.
Perbaikan: ubah \istilah{Collision Detection} menjadi \istilah{Continuous},
atau tebalkan lantainya. \\
\addlinespace
3 & \ind{c} &
Pertama, benda digerakkan melalui \code{transform.position} pada
\code{Update}, sehingga letaknya berubah di antara dua langkah fisika.
Kedua, \istilah{Interpolate} belum diaktifkan pada benda yang diikuti kamera.
Pemeriksaannya dengan melihat fungsi tempat benda digerakkan, dan dengan
memeriksa pengaturan \istilah{Rigidbody} pada \istilah{Inspector}. \\
\addlinespace
4 & \ind{b} &
\istilah{Is Trigger} harus ditandai, sekurang-kurangnya satu benda harus
memiliki \istilah{Rigidbody}, dan \istilah{layer} keduanya harus saling
diperiksa. Apabila salah satunya tidak dipenuhi, fungsinya tidak pernah
dipanggil sama sekali, tanpa pesan kesalahan apa pun. \\
\addlinespace
5 & \ind{c} &
Karena kamera harus dipindahkan menurut letak pemain yang sudah diperbarui
pada bingkai itu. \code{LateUpdate} dipanggil setelah seluruh \code{Update} dan
\code{FixedUpdate} selesai, sehingga letaknya sudah pasti. \\
\bottomrule
\end{tabularx}}
\end{table}
